% !TEX root = chapter-standalone.tex


\section{Ordering monotone functions}
\label{sec:ordering-monotone-functions}

Given \SY{preorders} $\posA$ and $\posB$, the set of \SY{monotone functions} $\posA \mto \posB$ can itself be made into a preorder: we compare two monotone functions ``point-wise'', by comparing their outputs for each input.

\begin{ctdefinition}[Pointwise order on monotone functions]
    \label{def:order-monotone-maps}
    \SYNDEF{order on monotone maps}
    Let $\posA$ and $\posB$ be \SY{preorders}, and let $\mapa, \mapb \colon \posA \mto \posB$ be \SY{monotone functions}.
    We say that $\mapa$ precedes $\mapb$ if, for every input, the output of $\mapa$ precedes the output of $\mapb$:
    \begin{equation}
        \prfdoubleperiod{
            \mapa \posleqof{\posA \mto \posB} \mapb
        }{
            \forall \posAel \setin \posAset \colon \mapa(\posAel) \posBleq \mapb(\posAel)
        }
    \end{equation}
    We write $\posB^\posA$ for the set of monotone functions $\posA \mto \posB$ equipped with this relation.
\end{ctdefinition}

The ``exponent'' notation $\posB^\posA$ is borrowed from sets, where $\setB^\setA$ denotes the set of all functions $\setA \sto \setB$ (\cref{sec:arithmetic-exponents-products}).

\begin{lemma}
    \label{lem:order-monotone-maps-preorder}
    Let $\posA$ and $\posB$ be \SY{preorders}.
    Then $\posB^\posA$ is a \SY{preorder}.
    If $\posB$ is a \SY{poset}, then $\posB^\posA$ is a \SY{poset}.
\end{lemma}
\begin{proof}
    Each property is checked one input at a time.
    Reflexivity: $\mapa(\posAel) \posBleq \mapa(\posAel)$ for all $\posAel$, by reflexivity of $\posBleq$.
    Transitivity: if $\mapa(\posAel) \posBleq \mapb(\posAel)$ and $\mapb(\posAel) \posBleq \mapc(\posAel)$ for all $\posAel$, then $\mapa(\posAel) \posBleq \mapc(\posAel)$ for all $\posAel$, by transitivity of $\posBleq$.
    Antisymmetry, when $\posB$ is a poset: if $\mapa(\posAel) \posBleq \mapb(\posAel)$ and $\mapb(\posAel) \posBleq \mapa(\posAel)$ for all $\posAel$, then $\mapa(\posAel) = \mapb(\posAel)$ for all $\posAel$ by antisymmetry of $\posBleq$, that is, $\mapa = \mapb$.
\end{proof}

Note that the order on $\posA$ plays no role in the definition of the order on $\posB^\posA$; it only determines which functions are monotone, and hence which elements $\posB^\posA$ has.

\begin{example}[Monotone functions into the booleans]
    \label{exa:bool-bool}
    There are four functions $\boolset \sto \boolset$, but only three of them are monotone functions $\Bool \mto \Bool$: the constant function with value $\false$, the identity, and the constant function with value $\true$.
    (Negation is not monotone, since $\false \posleqof{\Bool} \true$ but $\lnot \false = \true$ is not below $\lnot \true = \false$.)
    In $\Bool^\Bool$, these three functions form a chain:
    \begin{equation}
        (\text{constant } \false) \posleq \mapidat{\boolset} \posleq (\text{constant } \true).
    \end{equation}
    More generally, for any preorder $\posA$, monotone functions $\posA \mto \Bool$ correspond to \SY{upper sets} of $\posA$, and the point-wise order corresponds to inclusion of upper sets; we will use this in \cref{sec:dp-querying}.
\end{example}

\begin{example}[Rounding functions]
    \label{ex:rounding-functions}
    In this example we look at ``rounding functions'', which turn a real number into an integer.
    You might already know the ceiling function $\funceil$ (\cref{fig:ceil}) and the floor function $\funfloor$ (\cref{fig:floor}), which are defined as
    \begin{equation}\label{eq:funceil}
        \defmapcomma{
            \funceil
        }{
            \realswithleq
        }{
            \mto
        }{
            \tup{\wnumbers, \leq}
        }{
            \ela
        }{
            \min \makeset{\elb \setin \wnumbers \colon \elb \geq \ela }
        }
    \end{equation}
    and
    \begin{equation}\label{eq:funfloor}
        \defmapperiod{
            \funfloor
        }{
            \realswithleq
        }{
            \mto
        }{
            \tup{\wnumbers, \leq}
        }{
            \ela
        }{
            \max \makeset{\elb \setin \wnumbers \colon \elb \leq \ela }
        }
    \end{equation}
    The functions $\funceil$ and $\funfloor$ are monotone, since $\ela \Rleq \elc$ implies both $\funceil(\ela) \leq \funceil(\elc)$ and $\funfloor(\ela) \leq \funfloor(\elc)$.

    There exist many other rounding functions, commonly used by computers.
    For example, the function ``round to nearest, ties to even'' \cite{P754:2008:ISF} rounds a number to the closest integer, and in case of a tie it rounds to the even one (\cref{fig:rtntte}).
    For example, $3.2$ is mapped to $3$, $1.5$ is mapped to $2$, and $4.5$ is mapped to $4$.
    This is the formal definition:
    %
    \begin{widepar}
        \begin{equation}\label{eq:rtntte}
            \defmapperiod{
                \rtntte
            }{
                \realswithleq
            }{
                \mto
            }{
                \tup{\wnumbers, \leq}
            }{
                \ela
            }{
                \begin{cases}
                    \funfloor(\ela), & \ela - \funfloor(\ela) < \tfrac{1}{2},                                                   \\
                    \funceil(\ela),  & \ela - \funfloor(\ela) > \tfrac{1}{2},                                                   \\
                    \text{the even one of } \funfloor(\ela) \text{ and } \funceil(\ela), & \ela - \funfloor(\ela) = \tfrac{1}{2}
                \end{cases}
            }
        \end{equation}
    \end{widepar}
    It is monotone as well: as $\ela$ increases, the nearest integer can only increase, and the tie-breaking rule never makes it jump back.
    \begin{marginfigure}
        \includesag{floorceilhasse}
        \caption{}
        \label{fig:floorceilhasse}
        \todographics{\bernina: @Gioele: can we align id and rtntte in the center?}
    \end{marginfigure}

    To compare these functions with the identity, we view all of them as monotone functions $\realswithleq \mto \realswithleq$, using that every integer is a real number.
    Then, in $\realswithleq^{\realswithleq}$,
    \begin{equation}\label{eq:funfloor-mapid-funceil}
        \funfloor \posleq \mapidat{\reals} \posleq \funceil
    \end{equation}
    and
    \begin{equation}\label{eq:funfloor-rtntte-funceil}
        \funfloor \posleq \rtntte \posleq \funceil,
    \end{equation}
    while $\mapidat{\reals}$ and $\rtntte$ are not comparable: $\rtntte(0.4) = 0 < 0.4$, but $\rtntte(0.6) = 1 > 0.6$ (see \cref{fig:floorceilhasse}).
\end{example}

\begin{figure*}[b]
    \centering
    \subfloat[$\funceil$.
        \label{fig:ceil}]{
        \includesag{ceil}}
    \hfill
    \subfloat[$\funfloor$. \label{fig:floor}]{
        \includesag{floor}}
    \hfill
    \subfloat[$\rtntte$. \label{fig:rtntte}]{
        \includesag{rtntte}}
    \caption{Comparison of three rounding methods.}
\end{figure*}

\begin{example}[Comparing battery technologies]
    \label{exa:battery-technologies}
    A monotone function can describe how much of a resource a technology requires in order to provide a given functionality.
    Suppose that, for three battery technologies, the mass (in kg) needed to store the energy $\ela$ (in kWh) is
    \begin{equation}
        \mapa(\ela) = 1 + 4 \ela,
        \qquad
        \mapb(\ela) = 0.2 + 5 \ela,
        \qquad
        \mapc(\ela) = 0.2 + 4 \ela,
    \end{equation}
    where the constant term is the mass of the casing and electronics.
    All three are monotone functions $\tup{\nonNegReals, \leq} \mto \tup{\nonNegReals, \leq}$: storing more energy needs more mass.
    In the point-wise order, a smaller function describes a \emph{better} technology, since it needs less mass for every amount of energy.
    \begin{itemize}
        \item $\mapc \posleq \mapa$ and $\mapc \posleq \mapb$: the third technology is at least as good as the other two, for every amount of energy.
        \item $\mapa$ and $\mapb$ are not comparable: they are equal at $\ela = 0.8$, the second technology is lighter for smaller batteries ($\mapb(0.5) = 2.7 < 3 = \mapa(0.5)$), and the first is lighter for larger ones ($\mapa(2) = 9 < 10.2 = \mapb(2)$).
    \end{itemize}
    So whether $\mapa$ or $\mapb$ is preferable depends on how much energy we need to store; this is another instance of a trade-off.
\end{example}

\subsection{Composition and the order}

The point-wise order interacts well with composition of monotone functions.

\begin{lemma}
    \label{lem:composition-monotone-in-order}
    Let $\posA$, $\posB$, and $\posC$ be \SY{preorders}.
    \begin{enumerate}
        \item If $\mapa \posleqof{\posA \mto \posB} \mapa'$ and $\mapb \colon \posB \mto \posC$ is monotone, then $\mapa \mthen \mapb \posleqof{\posA \mto \posC} \mapa' \mthen \mapb$.
        \item If $\mapa \colon \posA \mto \posB$ is monotone and $\mapb \posleqof{\posB \mto \posC} \mapb'$, then $\mapa \mthen \mapb \posleqof{\posA \mto \posC} \mapa \mthen \mapb'$.
    \end{enumerate}
\end{lemma}
\begin{proof}
    \begin{enumerate}
        \item For every $\posAel$, we have $\mapa(\posAel) \posBleq \mapa'(\posAel)$, and since $\mapb$ is monotone, $\mapb(\mapa(\posAel)) \posCleq \mapb(\mapa'(\posAel))$.
        \item For every $\posAel$, the element $\mapa(\posAel)$ is an input of $\mapb$ and $\mapb'$, so $\mapb(\mapa(\posAel)) \posCleq \mapb'(\mapa(\posAel))$.
    \end{enumerate}
\end{proof}

Note that part 1 uses the monotonicity of $\mapb$, while part 2 uses no monotonicity at all.
Combining the two parts, if $\mapa \posleq \mapa'$ and $\mapb \posleq \mapb'$, then $\mapa \mthen \mapb \posleq \mapa' \mthen \mapb \posleq \mapa' \mthen \mapb'$.

\begin{example}[Improving a component]
    Continuing \cref{exa:battery-technologies}, suppose that the mass of the battery determines the price of the vehicle that carries it, through a monotone function $\stylemaps{p} \colon \tup{\nonNegReals, \leq} \mto \tup{\nonNegReals, \leq}$.
    Then the price of a vehicle able to store the energy $\ela$ is $\stylemaps{p}(\mapa(\ela))$ with the first technology and $\stylemaps{p}(\mapc(\ela))$ with the third.
    Since $\mapc \posleq \mapa$, part 1 of \cref{lem:composition-monotone-in-order} gives $\mapc \mthen \stylemaps{p} \posleq \mapa \mthen \stylemaps{p}$: replacing a component by a better one can only improve the system, provided that the rest of the system depends monotonically on it.
\end{example}
